\section{Diagnostic Design and Method Space}

SID papers now optimize different artifact properties: learned item indexing
and TIGER-style generative retrieval established the interface~\cite{hua2023indexids,rajput2023tiger},
while later work stresses collaborative alignment, end-to-end or
differentiable tokenization, representation-source sensitivity,
code-assignment diversity, collision qualification, variable-length
interfaces, drift, and deployment surfaces~\cite{zheng2023lcrec,zhu2024cost,wang2024letter,liu2024elit,fu2026diger,qiao2026textasvision,hu2026quasid,pan2026adasid,cheng2026capsid,baikalov2026staleness,ju2026snapchatsid}.
Table~\ref{tab:method-coverage} organizes this space by diagnostic facet rather
than by leaderboard rank.

\begin{table*}[t]
\caption{Facet coverage matrix. `M' denotes main PDF evidence, `A' artifact
repository evidence, `H' interface hook, and `--' no v0 evidence. The table is
coverage context, not a claim that all anchor methods are reproduced.}
\label{tab:method-coverage}
\footnotesize
\setlength{\tabcolsep}{2.4pt}
\begin{tabular}{@{}p{0.17\textwidth}p{0.30\textwidth}p{0.11\textwidth}ccccccc@{}}
\toprule
Facet & Anchors & v0 role & D1 & D2 & D3 & D4 & D5a & D6 & D7 \\
\midrule
A. Canonical residual SID & TIGER, GRID, RQ-KMeans~\cite{rajput2023tiger,ju2025grid} & GRID feature-text & M & M & M & M & M & -- & -- \\
B1. Collaborative / predictability & ReSID, LC-Rec, CoST, LETTER, DiscRec~\cite{liang2026resid,zheng2023lcrec,zhu2024cost,wang2024letter,liu2025discrec} & ReSID bounded & M & M & M & M & M & -- & -- \\
B2. Collision / capacity & QuaSID, AdaSID, CARD, CapsID~\cite{hu2026quasid,pan2026adasid,wei2026card,cheng2026capsid} & stressors/backlog & A & A & -- & A & A & -- & -- \\
B3. Ranking / retrieval & DIGER, ELIT, joint search-rec SID~\cite{fu2026diger,liu2024elit,penha2025jointsid} & interface hook & -- & -- & H & -- & H & -- & H \\
B4. Bottleneck / interface & AsymRec, CapsID, Long SID/ACERec~\cite{huang2026asymrec,cheng2026capsid,xia2026acerec} & future target & -- & -- & -- & H & H & -- & H \\
C. Drift / staleness & DACT, SID staleness~\cite{feng2026dact,baikalov2026staleness} & optional smoke & -- & -- & -- & -- & -- & A & -- \\
D. Deployment / search & Snapchat SID, SID-Coord, unified search-rec~\cite{ju2026snapchatsid,li2026sidcoord,penha2025jointsid} & motivation & -- & -- & -- & -- & H & -- & H \\
R0/control layer & RecList, Elliot, sanity SIDs~\cite{chia2022reclist,anelli2021elliot} & controls & A & A & A & A & A & -- & -- \\
\bottomrule
\end{tabular}
\end{table*}

\paragraph{D1--D4.}
D1 reports per-level usage, prefix counts, imbalance summaries, and dead or
underused code indicators. D2 measures duplicate full \sid{}s and prefix
collisions; it is a collision profile, not a causal harm estimate. A bounded
interaction-qualified controller is included because collision-aware work
argues that collisions are not uniformly harmful~\cite{hu2026quasid}. D3 asks
whether prefix neighborhoods recover item co-occurrence neighbors, keeping
metadata purity as auxiliary context rather than as collaborative quality. D4
splits items by popularity and measures whether full-code capacity and prefix
structure are allocated differently across head, mid, and tail items.

\paragraph{D5a--D7.}
D5a reports structural cost proxies: \sid length, prefix fan-out, duplicate
codes, and trie-like expansion pressure. These are not serving latencies; long
or variable SID methods make this distinction important~\cite{xia2026acerec,cheng2026capsid}.
D6 computes churn between tokenizer refreshes for drift-aware settings~\cite{feng2026dact,baikalov2026staleness}.
D7 is an interface hook for generator behavior, such as invalid paths,
next-token entropy, and duplicate generated candidates; it requires
\texttt{generator\_outputs} and is not claimed by the current evidence.
